% English translation of the author's revised Korean Introduction.
% The preceding Introduction is preserved unchanged in 01-introduction.tex.
\section{Introduction}\label{sec:introduction}

Vision-language models (VLMs) and vision-language-action (VLA) models connect
visual observations and language information to action selection or trajectory
prediction in autonomous driving. Recent work uses supervised fine-tuning and
reinforcement learning to improve language reasoning and action prediction
together~\cite{r1}. New proposals integrate future-scene modeling into planning
or maintain language-based memories of past observations to support decisions
over time~\cite{Wang2026DriveWLA,Yan2026ADMemo}. Evaluation also extends from
individual prediction accuracy to closed-loop driving and behavioral changes
under external instructions~\cite{Jia2026Bench2DriveVL,You2026Language2Drive}.

Despite these advances, reliably connecting a model's judgment to the required
behavior remains difficult. Drive the Thoughts reports only 74\% agreement between
predicted trajectories and reasoning traces judged
reliable~\cite{r18}. Language2Drive reports that models can succeed under one
instruction yet fail under a feasible alternative when the initial scene is
held fixed~\cite{You2026Language2Drive}. Generating an explanation or driving
successfully under one condition therefore does not establish the ability to
adjust behavior when external conditions change. We focus on using externally
supplied behavioral rules as conditions for action selection. The same scene
and candidates can admit different actions under different rules, so a model
must comply with the rule and select an action appropriate to that condition.

This requirement is distinct from explaining the reason for a driving
decision. Driving models that reason in language express how an observed
situation leads to a decision. Alpamayo-R1's Chain of Causation (CoC) is one
example: it learns causal relationships between observations and decisions
together with trajectory prediction~\cite{r1}. Such reasoning need not fully
specify an external rule's applicability conditions and numerical parameters
as a case-specific constraint. The scene and causal explanation can remain
valid while different speed bounds, waiting durations, or stop-release
conditions change the admissible actions. An action consistent with its
explanation can therefore violate a separately supplied rule.

Rule use requires both compliance and appropriate selection. Constraints
define the admissible actions, and the selection objective ranks actions
within that set; repeatedly choosing a waiting candidate admissible under
either rule may avoid violations while missing permitted progress. We
therefore measure violation rate and accuracy together, and rule-pair accuracy
requires both cases in a pair to be correct when only the rule changes.

Learning such selection raises the question of how training examples should
express the correspondence among scenes, applicable rules, and actions.
Providing a rule at inference makes the information available, but its effect
on selection requires assessment. Our preliminary Alpamayo-R1 diagnostics,
without additional training, also produced predictions violating explicitly
supplied speed and position-hold rules on sampled
PhysicalAI~\cite{nvidia2025physicalai} and converted
nuScenes~\cite{caesar2020nuscenes} inputs. In a separate controlled
candidate-selection study, we therefore test whether including constraints in
training improves compliance and rule-dependent selection when every model
receives the same rule at evaluation.

When constraints become learning information, their construction and
association with cases become data-quality concerns. Three difficulties arise.
First, one constraint combines information with different roles: entry after
pedestrian clearance involves an observed clearance time, a supplied waiting
duration, and a modeling assumption of no subsequent reoccupation. These roles
must be separated to trace the source of an error. Second, a fluent sentence
can be attached to the wrong case or encode a logical restriction different
from the intended one. Checks must therefore address correspondence with the
case, evidence, and rule as well as the sentence itself. Third, mechanical
checking requires a formal representation, whereas the model learns from
natural language. Both specification--sentence correspondence and the model's
behavioral use of the sentence require assessment.

Existing research provides foundations for this connection. nl2spec lets
users revise correspondences between language fragments and logical
subformulas~\cite{r9}. Semantic loss incorporates logical constraints into a
training objective to encourage constraint satisfaction~\cite{Xu2018}.
Instruction tuning shows that training with instructions can improve
performance on unseen task types~\cite{Wei2022}. Our preceding work showed
that natural-language guards appended to CoC improve action
selection~\cite{r4}. We connect these foundations in developing case-specific
behavioral constraints: binding constraints with numerical parameters to case
evidence, checking designated specification properties and sentence
correspondence, and supplying the resulting language alongside the original
reasoning as learning input.

We propose GuardSynth for this development process, addressing the three
difficulties in turn. It separates the roles of information by binding each
constraint field to an observation, a supplied rule, or a modeling assumption.
It makes the intended restriction checkable by representing applicability,
action restrictions, persistence and release, timing, and evidence in
Evidence-Bound Lifecycle Contracts (EBLC); supported specifications are lowered
to Satisfiability Modulo Theories (SMT) queries that check consistency and
designated properties within a restricted logical model. It connects the
checked meaning to learning by rendering the specification as Controlled
Natural Language (CNL) appended to the original CoC, with supported
fixed-grammar parsers recovering fields from the generated sentences to check
their correspondence with the specification.

The evaluation separates the behavioral effect of training with constraints
from specification checkability. In synthetic candidate-selection tasks with
Qwen3-VL-2B-Instruct~\cite{qwen2025qwen3vl2b}, fresh adapters share initialization
and evaluation inputs. Training with CNL under cross-entropy loss alone yields
three-seed mean violation rates of at most 1.56\% and rule-pair accuracy of at
least 86.98\% across pedestrian entry and the two position-hold representations,
improving on CoC-only training. Across six behavioral tasks, both natural-language
constraints and EBLC-derived CNL improve rule-dependent selection. These
findings support including case-specific constraints in training data as well
as supplying them at inference.

A separate fault-injection evaluation of the two learning specifications
detects all 288 errors of six supported types and accepts all 48 valid cases.
It misses 48 cases in which the evidence record and specification are wrong
together, identifying a boundary between internal correspondence checks and
independent evidence confirmation. Source- and sentence-correspondence checks
were added after training for this evaluation; their results are interpreted
separately from the causes of behavioral improvement.

Our contributions are:
\begin{enumerate}
\item A controlled evaluation of training with case-specific constraints.
Identical initialization and evaluation inputs isolate the training-text
comparison, while violation rate, accuracy, and rule-pair accuracy distinguish
compliance from rule-dependent selection.
\item A GuardSynth development path from evidence binding to natural-language
learning inputs. It separates observations, supplied rules, and assumptions,
constructs EBLC specifications, and connects restricted logical checking and
field--sentence correspondence checks to CNL appended to CoC.
\item An evaluation of checking layers and their detection scope for learning
specifications. Fault injection into the two specifications used to generate
learning inputs identifies complementary checks and their limits, separately
from model compliance.
\end{enumerate}

Section~\ref{sec:related-work} reviews related work, and
Section~\ref{sec:method} presents the problem setting and GuardSynth.
Sections~\ref{sec:experimental-setup} and~\ref{sec:results} describe the
experimental setup and results. Section~\ref{sec:discussion} discusses
implications and limitations, and Section~\ref{sec:conclusion} concludes.
